% Supporting TeX fragment, intended for the existing raw-payoff proof.
% This contains a checked derivation, not a claim of theoretical priority.
\section{A smoothing-first arithmetic certificate}
\label{sec:raw-smoothing-first}

The complete pointwise graph certificate remains necessary for unsigned
domains, scalar input domains and the absence of modular overflow. Its
pointwise output error need not be the sharpest certificate of the prepared
mean. I now derive a second bound for the same emitted arithmetic and take
the smaller of the two bounds. This changes the financial certification,
prior enlargement and estimator allowance; it does not change the circuit.
The following argument separates rounding which depends on the inner
normal label from rounding which depends only on the outer coordinate.
Consequently, it does not differentiate a finite atomic law or a payoff
across its kinks.

Write $s=2^q$ and $u_q=s^{-1}$. The nearest dyadic rounding convention is
\[
 \mathcal R_q(v)=s^{-1}\lfloor sv+1/2\rfloor.
\]
All quantities in this section are real values decoded from the guarded
integer graph, rather than integers interpreted modulo a register width.
Let $y$ be any realised outer-grid coordinate. The exact targets are
$m=m_G(y)$ and $\Delta=\log T(y)\geq0$. The shifted affine node has ideal
value $M=m+B_G$, where $B_G$ is the exact offset inherited from the retained
graph. Its implemented value is $\widehat M$. The implemented nonnegative
logarithm node is $\widehat\Delta$. Let the retained node certificate give
the uniform bounds
\[
 |\widehat M-M|\leq e_m,
 \qquad |\widehat\Delta-\Delta|\leq e_\Delta,
 \qquad \widehat\Delta\geq0.
\]
The words ``realised outer-grid coordinate'' are essential: Gaussian
outer-node rounding is already included in the inherited financial
transfer. The arithmetic certificate treats those realised coordinates
as its exact inputs and does not charge their rounding a second time.

Put
\[
 C=J-B_G,\quad O=J+\lambda L,\quad
 \widehat C=\mathcal R_q(C),\quad
 \widehat O=\mathcal R_q(O),
\]
and define $\epsilon_C=\widehat C-C$ and
$\epsilon_O=\widehat O-O$. Their certified absolute bounds are
$e_C$ and $e_O$. These constants are independent of the inner normal
label. In this graph the lower-endpoint additions and the upper-endpoint
addition are exact integer additions. Their real decoded values therefore
agree with ordinary real addition, once the retained whole-graph overflow
and domain guards pass. No additional addition-rounding term is present.

Let $(I_\lambda,d_\lambda)$ be exactly the coefficient pair returned by
the retained fixed-point scaling emitter and set
\[
 \widehat\lambda=I_\lambda 2^{-d_\lambda},
 \qquad \delta_\lambda=\widehat\lambda-\lambda,
 \qquad
 e_{\rm floor}=\begin{cases}u_q,&d_\lambda>0,\\0,&d_\lambda=0.\end{cases}
\]
For the realised quantile word $v=\widetilde N+L\in[0,2L]$, the scale
node decodes to
\[
 \mathsf S_\lambda(v)
 =s^{-1}\left\lfloor (sv)I_\lambda 2^{-d_\lambda}\right\rfloor
 =\widehat\lambda v+\rho(v),
 \qquad |\rho(v)|\leq e_{\rm floor}.
\]
When $d_\lambda=0$, the product is an integer and $\rho(v)=0$ exactly.
The inequality with $u_q$ is outward and remains valid at dyadic endpoints.
Thus, with
\[
 E_s=2L|\delta_\lambda|+e_{\rm floor},\qquad
 \widehat m=\widehat M-B_G+\epsilon_C-\epsilon_O,
\]
the two conceptual signed log arguments before positive subtraction are
exactly
\begin{align*}
 \widehat X_G
 &=\widehat M+\widehat C+
     \mathsf S_\lambda(\widetilde N+L)-\widehat O
   =\widehat m+\lambda\widetilde N+\zeta(\widetilde N),\\
 \widehat X_A
 &=\widehat X_G+\widehat\Delta,
 \qquad
 \zeta(n)=\delta_\lambda(n+L)+\rho(n+L),
 \qquad |\zeta(n)|\leq E_s.
\end{align*}
The signed expressions are mathematical identities used to identify the
target. The circuit itself computes their positive parts by guarded
positive subtraction; it does not store an unguarded signed logarithm.
In particular,
\[
 |\widehat m-m|\leq E_m:=e_m+e_C+e_O.
\]

\begin{lemma}[Scalar implementation error after exact payoff clamps]
\label{lem:raw-payoff-scalar-error}
Let $b=1/64$ and let the implemented exponential map $\widehat E$ satisfy
\[
 \sup_{t\in[0,b]\cap s^{-1}\mathbb Z}
 |\widehat E(t)-e^t|\leq e_{\exp}.
\]
With $K=100$ and $H=1$, the graph payoff map obeys
\[
 |\widehat f(x)-f(x)|\leq100e_{\exp}
\]
for every guarded decoded log argument $x$.
\end{lemma}
\begin{proof}
The positive-part map and the clamps at $b$ and one are exact and
$1$-Lipschitz. Both constants $b$ and one are exactly representable at
$q=28$. The scalar graph's second input clamp at $b$ is idempotent and
exact. The operation which subtracts one positively is exact because
the constant one is exact and the positive-subtraction guards pass.
Scaling by the integer $100$ has $d=0$, so it is an exact integer product
without a fixed-point floor. The difference between implemented and exact
exponentials is therefore multiplied by at most $100$ and is not enlarged
by either following clamp. The clamp decomposition established earlier
identifies the exact composed payoff with $f(x)$, completing the proof.
\end{proof}

The scalar certificate used in this lemma is its
\texttt{output\_error\_upper}, comprising nodal rounding, smooth
interpolation and the interpolation floor. It excludes the propagated
argument-error charge. Argument perturbations are treated separately
below. This distinction is required to avoid both omitting a perturbation
and charging it twice. Allow distinct bounds $e_{\exp,G}$ and
$e_{\exp,A}$, although the actual graph uses the same table and hence the
same bound on both sides.

\begin{theorem}[Smoothing-first conditional mean error]
\label{thm:raw-smoothing-first}
Assume the whole-graph domain and overflow certificate passes, the
quantile CDF bound is $D_N$, the shifted-node and coefficient bounds above
hold, and $\nu_0\geq\sup_x\phi(x)$. Let $p(y)$ be the exact probability
of the prepared good flag conditional on the realised outer coordinate.
Then
\begin{align*}
 |p(y)-R_1(y)|
 &\leq\eta_{\rm smooth},\\
 \eta_{\rm smooth}
 &=\min\left\{1,
     100(e_{\exp,G}+e_{\exp,A})+202E_s+2D_N
     +\frac{\nu_0}{\lambda}(2E_m+e_\Delta)\right\}.
\end{align*}
If $\eta_{\rm point}$ is the independently valid pointwise graph/CDF
bound, the smaller bound
$\eta=\min\{\eta_{\rm point},\eta_{\rm smooth}\}$ is valid uniformly
over the same realised outer grid.
\end{theorem}
\begin{proof}
The payoff is absolutely continuous, with its almost-everywhere derivative
\[
 f'(x)=100e^x\mathbf1_{(0,\ell)}(x),
 \qquad \ell=\log(101/100).
\]
Thus $f$ is globally $101$-Lipschitz. Replacing
$\widehat X_G$ and $\widehat X_A$ by
$\widehat m+\lambda\widetilde N$ and
$\widehat m+\widehat\Delta+\lambda\widetilde N$ changes the difference
of exact payoffs by at most $2\cdot101E_s=202E_s$ for each normal label.
The scalar lemma changes the difference by at most
$100(e_{\exp,G}+e_{\exp,A})$. The final positive subtraction and clamp
are the metric projection onto $[0,1]$; that projection is $1$-Lipschitz
and fixes the exact payoff difference because
$\widehat\Delta\geq0$. They cannot enlarge this bound.

For completeness, the represented output is in fact the unprojected
difference of represented payoffs. The upper represented endpoint is
the lower endpoint plus the nonnegative word $\widehat\Delta$, without
overflow. Each preceding clamp and positive subtraction is monotone.
The rounded exponential node integers are nondecreasing. Within a cell,
the interpolator returns its base plus the floor of a nonnegative
increment. It is nondecreasing and cannot exceed the next node; at the
next cell it begins at that node. Hence the interpolator is monotone on
the complete dyadic input domain. Positive scaling and the final payoff
clamp preserve order, giving $0\leq\widehat f(\widehat X_G)
\leq\widehat f(\widehat X_A)\leq1$. Exact final subtraction therefore
produces their difference. The uniform comparator implements that
dyadic difference as a flag probability exactly.

After these pointwise replacements, the finite-inner-law target is
$\mathbb E h_{\widehat m,\widehat\Delta}(\widetilde N)$.
Both $\widehat m$ and $\widehat\Delta$ depend only on outer labels,
constants and outer arithmetic. They are independent of the inner normal
label, as is required here. The existing bounded-variation CDF lemma
applies at these represented endpoints, including atomic or duplicate
quantiles, and bounds its difference from
$\mathbb E h_{\widehat m,\widehat\Delta}(N)$ by $2D_N$.
Normal-node rounding is entirely charged by $D_N$ and does not enter
$E_s$ or $E_m$.

It remains to compare two continuous conditional means. Define
\[
 P_\lambda(t)=\mathbb E f(t+\lambda N).
\]
The nonnegative derivative $f'$ has mass
$\int f'(x)\,dx=100(e^\ell-1)=1$. Convolution with the Gaussian density
therefore gives
\[
 P_\lambda'(t)
 =\int_{\mathbb R} f'(x)\frac1\lambda
             \phi\left(\frac{x-t}{\lambda}\right)\,dx,
 \qquad 0\leq P_\lambda'(t)\leq\frac{\nu_0}{\lambda}.
\]
One justification of the derivative identity is to write the difference
quotient of $f(t+\lambda N)$ as an integral of its bounded almost-everywhere
derivative and use dominated convergence. Changing variables yields the
displayed convolution. Alternatively, the finite positive derivative
measure $df$ can be convolved with the smooth Gaussian. Either argument
integrates the continuous normal before applying the slope bound and
never differentiates the finite atomic normal approximation.

The Lipschitz bound for $P_\lambda$ now yields
\begin{align*}
 &\left|P_\lambda(\widehat m+\widehat\Delta)
          -P_\lambda(\widehat m)
          -P_\lambda(m+\Delta)+P_\lambda(m)\right|\\
 &\qquad\leq\frac{\nu_0}{\lambda}
        \left(|\widehat m-m+\widehat\Delta-\Delta|
              +|\widehat m-m|\right)
 \leq\frac{\nu_0}{\lambda}(2E_m+e_\Delta).
\end{align*}
Adding the three independently justified errors proves the untruncated
bound. Both the prepared probability and $R_1(y)$ lie in $[0,1]$, so
truncating the error at one is valid. Taking the smaller of two valid
uniform upper bounds is valid as well.
\end{proof}

The analytic financial coefficient, outer finite law, its external mass
and continuous financial transfer remain unchanged. The complete monetary
certificate is rebuilt with this new $\eta$, and the same ideal finite-mean
prior is enlarged by this new $\eta$. The phase-estimation allowance,
coverage plan and full logical resource charge must be rebuilt afterwards.
Neither a new estimator cost nor a strengthened comparative conclusion
follows merely from deriving a smaller arithmetic bound. This is a
project-specific derivation from established smoothing, monotonicity,
fixed-point rounding and CDF transfer principles; theoretical priority is
not asserted.
